\chapter{A Brief Prehistory of Doxonomics}
\label{ch:prehistory}

\epigraph{The philosophers have only interpreted the world, in various ways; the point is to change it.}{Karl Marx, \textit{Theses on Feuerbach} (1845)}

\noindent
Doxonomics is new as a named discipline.
But the insight that beliefs are shaped by material power is ancient.
This chapter traces the intellectual lineage, extracting from each predecessor the formal insight that doxonomics inherits.
The genealogy is not meant to be exhaustive. It is organized around the question: \emph{who saw the conversion problem, and what tools did they develop to analyze it?}


\section{Sophists and the Market for Persuasion}
\label{sec:ch3-sophists}

The Greek sophists were, in effect, the first consultants in belief production.
Protagoras, Gorgias, Thrasymachus, and their colleagues sold rhetorical skill, the capacity to make the weaker argument appear stronger.
They charged fees.
They competed for students and patrons.
They were, in modern terms, a market for persuasion services.

Plato attacked them for it, famously in the \textit{Gorgias} and the \textit{Republic}.
His objection was that rhetoric separated persuasion from truth: the sophist could make any proposition plausible, regardless of its accuracy, provided the audience was susceptible and the speaker was skilled.
This is already a proto-doxonomic critique.

\begin{remark}
The sophists reveal the first doxonomic principle: \emph{the ability to shape belief is an economic good with a market price}.
In modern notation, there exists a market-clearing price $p^*$ for persuasion services, which implies that wealthier actors can purchase more belief-shaping capacity, a precursor to the concept of epistemic capital ($\ecap$).
The Platonic response (that truth should govern belief) is the normative position that doxonomics, as a descriptive science, neither endorses nor rejects but takes as a benchmark against which to measure distortion.
\end{remark}

What the sophists lacked was a structural analysis.
They understood individual persuasion.
They did not theorize the institutional systems by which persuasion becomes common sense across an entire polity.


\section{Hobbes, Mandeville, and the Construction of Selfish Anthropology}
\label{sec:ch3-hobbes}

\textcite{hobbes1651} proposed that the natural condition of humanity is a war of all against all, requiring absolute sovereignty to impose order.
\textcite{mandeville1714} argued in \textit{The Fable of the Bees} that private vices produce public benefits. Selfishness is not only natural but socially productive, and attempts to enforce virtue would impoverish society.

These are not merely philosophical positions.
They are \emph{anthropological regimes} (Definition~\ref{def:anthropological-regime}) whose institutional installation required centuries of work: in political theory, in economics curricula, in management doctrine, in legal reasoning.
By the late twentieth century, the Hobbesian-Mandevillean anthropology had become operational common sense in most Western institutions, not because it won a philosophical argument, but because it was embedded in the design of markets, contracts, incentive systems, and governance structures.

\begin{remark}
\label{rem:hobbes-doxonomic}
The doxonomic question about Hobbes is not ``was he right?'' but ``what institutional investment was required to make his anthropology the default assumption of modern governance, and what alternative anthropologies were crowded out in the process?''
\textcite{bowles2011} and \textcite{ostrom1990} provide evidence that the crowded-out alternatives had substantial empirical support.
\end{remark}

\textcite{smith1776} occupies an ambiguous position.
The \textit{Wealth of Nations} is routinely cited as a founding text of the self-interest paradigm (``It is not from the benevolence of the butcher, the brewer, or the baker that we expect our dinner'').
But Smith also wrote \textit{The Theory of Moral Sentiments} (1759), which grounds human sociality in sympathy, fellow-feeling, and the desire for mutual recognition.
The selective amplification of Smith-as-self-interest-theorist, and the suppression of Smith-as-moral-philosopher, is itself a doxonomic achievement, one that \textcite{mcCloskey2006} has documented in detail.


\section{Marx: Ideology as Material Production}
\label{sec:ch3-marx}

\textcite{marx1846} proposed the most radical version of the thesis: the ruling ideas of any epoch are the ideas of the ruling class, produced not by the force of argument alone but by control over the means of mental production.
In the same way that the class which owns the factories controls the production of goods, it also controls, through ownership of newspapers, publishing houses, universities, and churches, the production of ideas.

\begin{principle}[Marxian Doxonomic Principle]
\label{prin:marx}
\index{Marx}
The class that controls the means of material production simultaneously controls the means of intellectual production.
In doxonomic terms: the funding vector $\fundingvec$ is not exogenous but is determined by the distribution of economic power.
Formally, let $W$ be the wealth distribution and $\fundingvec = g(W)$ for some function $g$ that maps material power to doxonomic investment.
If $g$ is monotonically increasing (wealthier actors invest more in belief production), then doxonomic output is structurally biased toward the interests of the wealthy.
\end{principle}

Marx's contribution to doxonomics is threefold:
\begin{enumerate}[nosep]
  \item the recognition that ideas have material conditions of production,
  \item the claim that the content of dominant ideas systematically reflects the interests of dominant groups,
  \item the concept of ideology as \emph{false consciousness}: beliefs that serve the interests of the powerful while being experienced by the powerless as their own convictions.
\end{enumerate}

His limitation, from a doxonomic standpoint, was insufficient formalization and insufficient attention to the mechanisms of belief transmission.
Marx knew \emph{that} ideas served class interests; he did not provide the accounting tools to trace \emph{how much} was spent, \emph{through which institutions}, with \emph{what measurable effects}.
He also underestimated the autonomy of institutional processing: ideas are not simply piped from the ruling class to the population but are transformed, filtered, and repackaged by intermediary institutions with their own logics.

\textcite{althusser1970} partially addressed this with the concept of \emph{ideological state apparatuses}\index{ideological state apparatus}: schools, churches, media, family, and legal systems that reproduce ideology not through direct coercion but through daily institutional practice.
In doxonomic terms, Althusser identified the institutional ecology $\insteco$ as a set of belief-processing functions that operate with relative autonomy from direct class instruction.


\section{Gramsci: Hegemony}
\label{sec:ch3-gramsci}

\textcite{gramsci1971} introduced \emph{hegemony}\index{hegemony}: the process by which a ruling group secures consent not through coercion but through cultural and intellectual leadership.
Hegemony operates through civil society (schools, churches, media, professional associations, cultural institutions) rather than through the state alone.

Gramsci's key innovation was the recognition that domination through belief is \emph{more effective and more durable} than domination through force.
A state that relies on police is fragile; a state whose citizens genuinely believe that the existing order is natural, just, or inevitable barely needs police at all.

In doxonomic terms, Gramsci's contribution is the recognition that $\insteco$ (the institutional ecology) is the primary site of belief production, and that durable power requires not just force but narrative infrastructure $\narr$.
He also introduced a crucial dynamic concept: hegemony is never total, always contested, and must be continuously \emph{renewed} through institutional work.
This anticipates the doxonomic concept of maintenance cost: even a captured belief requires ongoing investment to remain dominant.

\begin{remark}
Gramsci also theorized \emph{counter-hegemony}: the construction of alternative common sense by subordinate groups through their own institutional networks (workers' councils, alternative media, pedagogical movements).
This prefigures the counter-doxonomic strategies of Chapter~\ref{ch:counter-doxonomic}.
\end{remark}


\section{Lippmann and Dewey: The Public Opinion Debate}
\label{sec:ch3-lippmann-dewey}

In the 1920s, two American thinkers framed a debate that remains central to doxonomics.

\textcite{lippmann1922} argued in \textit{Public Opinion} that ordinary citizens are incapable of understanding the complexity of modern governance.
They perceive the world through ``stereotypes,'' simplified pictures in their heads, produced by media and institutional sources rather than by direct experience.
Public opinion is therefore not the spontaneous expression of citizen judgment but a \emph{manufactured product}, shaped by those who control the information environment.

Lippmann's doxonomic insight: the production function $\Phi$ is controlled by a relatively small number of institutional actors (editors, publishers, experts, officials), and the population's beliefs are the \emph{output} of this function, not its independent input.

\textcite{dewey1927} responded in \textit{The Public and Its Problems} that the solution was not technocratic management of opinion (Lippmann's implicit conclusion) but the \emph{democratization of communication}.
If the public is poorly informed, the remedy is not to hand belief production to experts but to create institutional conditions (free media, public education, open deliberation) under which citizens can form genuine judgments.

\begin{remark}
The Lippmann-Dewey debate maps directly onto a tension within doxonomics.
Lippmann's position leads to a \emph{diagnostic} doxonomics: mapping how belief is produced by elites and processed by publics.
Dewey's position leads to a \emph{constructive} doxonomics: designing institutional environments for democratic belief formation.
This book attempts both, but the tension is real and unresolved.
Chapter~\ref{ch:epistemic-democracy} returns to it.
\end{remark}


\section{Bernays: Engineering Consent}
\label{sec:ch3-bernays}

\textcite{bernays1928} made the conversion problem explicit.
Drawing on his uncle Freud's psychology, Bernays developed techniques for converting corporate funding into public opinion: the staged event, the third-party endorsement, the appeal to unconscious desire, the manipulation of group psychology.

His 1929 ``Torches of Freedom'' campaign (in which he hired women to smoke cigarettes in the New York Easter Parade, framing smoking as a feminist act) is a canonical example of doxonomic engineering.
The client was the American Tobacco Company.
The narrative output was ``smoking = women's liberation.''
The channel was news media (journalists covered the stunt as a news event, lending it third-party credibility).
The behavioral consequence was a measurable increase in women's smoking.

\begin{remark}
Bernays is the first explicit \emph{doxonomic engineer}: a practitioner who consciously understood and optimized the production function $\Phi: \R^m_{\geq 0} \times \insteco \to \narr$.
Public relations, as he conceived it, is applied doxonomics, the deliberate, funded transformation of institutional inputs into belief outputs.
His candor was remarkable.
The opening of \textit{Propaganda}: ``The conscious and intelligent manipulation of the organized habits and opinions of the masses is an important element in democratic society.
Those who manipulate this unseen mechanism of society constitute an invisible government which is the true ruling power of our country.''
\end{remark}

\textcite{ellul1965} extended this analysis.
Where Bernays focused on technique, Ellul argued that propaganda is not merely a tool used by states or corporations but a \emph{total social phenomenon}: modern societies produce propaganda as naturally as they produce consumer goods, because the complexity of modern life creates a demand for simplified interpretive frameworks that only organized narrative production can supply.
In doxonomic terms, Ellul's point is that the demand side of the conversion problem is endogenous: complex societies create the audience conditions under which doxonomic production becomes effective.


\section{The Frankfurt School: Culture Industry}
\label{sec:ch3-frankfurt}

\textcite{adorno1944} introduced the concept of the \emph{culture industry}\index{culture industry}: the industrial-scale production of cultural goods (film, music, radio, magazines) that standardize consciousness and integrate populations into the existing order.

The Frankfurt School's doxonomic insight was that \emph{entertainment is a channel of belief production}, not merely a diversion from it.
A Hollywood film that represents success as individual achievement, romance as consumption, and conflict as resolved through violence is producing an anthropological regime, not by arguing for it but by \emph{narrating} it as the natural shape of human life.

In doxonomic terms, the culture industry operates at layers 3--5 of the layered ontology (Principle~\ref{prin:layered-ontology}): practical schemas, identity commitments, and affective orientations.
It shapes not what people explicitly believe but what they desire, expect, fear, and find normal.
This makes it both more powerful and harder to trace than propositional propaganda.

The limitation of the Frankfurt School, from a doxonomic standpoint, was a tendency toward totalization: the culture industry was presented as so pervasive that resistance seemed impossible.
Later work (by Stuart Hall, the Birmingham cultural studies tradition, and audience reception theorists) showed that audiences are not passive consumers but active interpreters who negotiate, resist, and reinterpret cultural products \autocite{hall1980}.


\section{Foucault: Discourse, Power, and the Production of Truth}
\label{sec:ch3-foucault}

\textcite{foucault1971} shifted attention from who speaks to the \emph{rules of discourse}: the conditions under which a statement counts as knowledge, who is authorized to make it, and what institutional apparatus validates it.

In Foucault's analysis, power does not merely suppress truth; it \emph{produces} truth.
Each society has its ``regime of truth'': the types of discourse it accepts as true, the mechanisms that enable one to distinguish true from false statements, the techniques for acquiring truth, and the status of those who are charged with saying what counts as true.

In doxonomic terms, Foucault's contribution is the insight that $\credibility{k}$ is not a fixed property of a channel but is itself socially produced.
The rules determining what counts as expertise, evidence, or common sense are themselves objects of doxonomic investment.
This is the fourth dimension of doxonomic power (Definition~\ref{def:doxonomic-power} in Chapter~\ref{ch:doxonomic-power}): the power to shape the meta-rules of belief evaluation.

\begin{remark}
Foucault pushes doxonomics toward a second-order analysis.
The first-order model asks: how are beliefs produced?
The Foucauldian question asks: how are the \emph{rules for evaluating beliefs} produced?
Who decides what counts as evidence, what counts as expertise, what counts as a reasonable question?
These meta-rules are themselves doxonomic products, and among the most consequential ones.
A society that controls the definition of ``science,'' ``expert,'' and ``evidence'' controls the credibility weights $\credibility{k}$ for all channels simultaneously.
\end{remark}


\section{Bourdieu: Symbolic Power, Doxa, and Reproduction}
\label{sec:ch3-bourdieu}

\textcite{bourdieu1984, bourdieu1991} provided the richest pre-doxonomic framework.
His contributions map onto doxonomic concepts with unusual precision:

\begin{center}
\begin{tabular}{ll}
\toprule
\textbf{Bourdieu} & \textbf{Doxonomics} \\
\midrule
Symbolic capital & Epistemic capital $\ecap$ \\
Doxa & Common-sense capture (Definition~\ref{def:doxa}) \\
Field & Institutional ecology $\insteco$ \\
Habitus & Practical schema (layer 3) \\
Cultural reproduction & Belief production chain $\bpc$ \\
Misrecognition & Invisible as a claim (doxa condition 4) \\
\bottomrule
\end{tabular}
\end{center}

Bourdieu's analysis of the educational system as a reproduction mechanism is particularly relevant.
In \textit{Distinction}, he showed that the French educational system does not merely transmit knowledge; it converts inherited economic capital into credentialed cultural capital, laundering class privilege through the appearance of meritocratic achievement.
The school system is simultaneously an institution that produces knowledge and an institution that produces the \emph{belief} that the knowledge hierarchy reflects natural differences in ability.

What Bourdieu lacked was the accounting machinery.
He showed that symbolic power reproduces itself through institutions but did not build the input-output tables, trace the funding flows, or estimate the conversion rates.
Doxonomics inherits his sociology and adds the ledger.


\section{Herman and Chomsky: The Propaganda Model}
\label{sec:ch3-herman-chomsky}

\textcite{herman1988} provided the most systematic media-level doxonomic model.
The propaganda model identifies five ``filters'' through which media content passes before reaching the public:

\begin{enumerate}[nosep]
  \item \textbf{Ownership}: media corporations are owned by large conglomerates with interests in maintaining the status quo.
  \item \textbf{Advertising}: the primary revenue source for most media, giving advertisers structural influence over content.
  \item \textbf{Sourcing}: dependence on official sources (government, corporate PR) for raw material.
  \item \textbf{Flak}: organized campaigns to discredit or punish media that deviate from acceptable narratives.
  \item \textbf{Ideology}: shared assumptions (anticommunism in the original formulation; market fundamentalism in later versions) that frame what counts as newsworthy.
\end{enumerate}

In doxonomic terms, the five filters are structural parameters of the narrative infrastructure $\narr$ that systematically bias output without requiring any explicit conspiracy.
The propaganda model's great strength is that it is \emph{testable}: it generates predictions about which stories will be covered, how they will be framed, and which perspectives will be excluded, based on structural properties of the media system rather than on the intentions of individual journalists.

Its limitation is its focus on media.
Doxonomics extends the same structural logic to the full institutional landscape: education, think tanks, platforms, professional training, entertainment, and the built environment.


\section{Colonial Knowledge Production}
\label{sec:ch3-colonial}

The most large-scale doxonomic operations in history were not conducted by think tanks or advertising agencies.
They were conducted by colonial empires, which invested enormous resources in producing knowledge systems (racial taxonomies, civilizational hierarchies, linguistic classifications, legal codes) that justified extraction and governance.

\textcite{said1978} showed that ``Orientalism'' was not merely prejudice but an institutionally sustained body of knowledge: endowed chairs at universities, specialized journals, government translation offices, and training programs for colonial administrators, all dedicated to producing authoritative accounts of the colonized world.
The narrative infrastructure was vast: the British Indian Civil Service alone trained thousands of administrators in a specific interpretive framework for understanding ``the Oriental mind.''

\textcite{fanon1961} analyzed the internalization stage: how colonial education systems replaced indigenous knowledge with the colonizer's worldview, producing subjects who internalized their own subordination.
The colonial school taught not just the colonizer's language but the colonizer's anthropology: the belief that European civilization was the pinnacle of human achievement, that indigenous cultures were backward, and that colonial rule was a civilizing mission.

In doxonomic terms, colonialism built the most extensive narrative infrastructure $\narr$ in pre-digital history, reaching hundreds of millions through missionary schools, colonial administration, and restructured legal systems.
The persistence of these narratives after formal decolonization (what decolonial theorists call \emph{coloniality}\index{coloniality}) is a doxonomic phenomenon: the institutional infrastructure outlasts the political regime that built it.
Chapter~\ref{ch:race-coloniality-gender} develops this analysis.


\section{The Digital Turn: Platforms as Doxonomic Infrastructure}
\label{sec:ch3-digital}

The early twenty-first century introduced a qualitative shift in doxonomic infrastructure.
Social media platforms (Facebook, YouTube, Twitter/X, TikTok) became simultaneously:
\begin{itemize}[nosep]
  \item \emph{channels} of narrative circulation (reaching billions),
  \item \emph{institutions} of narrative processing (algorithmic curation decides what is amplified),
  \item \emph{data-collection systems} that enable micro-targeted doxonomic operations (personalized belief exposure based on behavioral profiles).
\end{itemize}

\textcite{zuboff2019} analyzed this as ``surveillance capitalism'': the extraction and commodification of behavioral data for predictive and manipulative purposes.
In doxonomic terms, platforms represent a new kind of narrative infrastructure in which the exposure matrix $E$ (who sees what) is not determined by editorial judgment or public policy but by engagement-maximizing algorithms whose objective function is revenue, not truth or public welfare.

This makes the conversion problem both easier and harder to study.
Easier, because platform advertising data provides unprecedented visibility into doxonomic spending (Meta's Ad Library, Google's Transparency Report).
Harder, because algorithmic curation operates at a speed and granularity that traditional doxonomic accounting cannot yet match.


\section{Synthesis: What Doxonomics Inherits}
\label{sec:ch3-synthesis}

Each predecessor contributed a piece.
Doxonomics assembles them:

\begin{center}
\begin{tabular}{lp{8cm}}
\toprule
\textbf{Predecessor} & \textbf{Doxonomic Inheritance} \\
\midrule
Sophists & Persuasion is a priced good; $\ecap$ has a market. \\
Hobbes/Mandeville & Anthropological regimes are installed, not discovered. \\
Marx & $\fundingvec$ is endogenous to the wealth distribution. \\
Althusser & $\insteco$ operates with relative autonomy from direct class instruction. \\
Gramsci & Durable power requires $\narr$, not just coercion; hegemony must be continuously renewed. \\
Lippmann & $\Phi$ is controlled by a small number of institutional actors. \\
Dewey & The remedy is democratized communication, not technocratic management. \\
Bernays & Applied doxonomics: optimizing $\Phi$ for commercial and political clients. \\
Ellul & Demand for simplified belief is endogenous to social complexity. \\
Frankfurt School & Entertainment is a doxonomic channel operating at layers 3--5. \\
Foucault & $\credibility{k}$ is itself socially produced (second-order doxonomics). \\
Bourdieu & Symbolic capital, doxa, reproduction: the richest pre-doxonomic vocabulary. \\
Herman/Chomsky & Structural filters produce systematic bias without conspiracy. \\
Said/Fanon & Colonial empires built the largest pre-digital $\narr$ in history. \\
Zuboff & Platforms as algorithmic doxonomic infrastructure with micro-targeting. \\
\bottomrule
\end{tabular}
\end{center}

What none of these predecessors provided, and what doxonomics aims to build, is a \emph{unified formal framework} that integrates accounting, network analysis, causal inference, measurement, and dynamical modeling into a single disciplinary structure.
The tools exist in scattered form across economics, sociology, media studies, and computer science.
The task of this book is to assemble them.


\section{A Timeline}

\begin{center}
\begin{tikzpicture}[
  every node/.style={font=\footnotesize},
  name/.style={font=\footnotesize\bfseries, rotate=45, anchor=north east},
]
  \draw[thick, -Stealth] (0,0) -- (15,0);

  \foreach \x/\t/\d in {
    0.3/Hobbes/1651,
    1.5/Mandeville/1714,
    2.7/Marx/1846,
    3.9/Lippmann/1922,
    5.0/Bernays/1928,
    6.1/Adorno/1944,
    7.3/Ellul/1962,
    8.5/Gramsci/1971,
    9.5/Foucault/1971,
    10.5/Said/1978,
    11.5/Bourdieu/1984,
    12.5/Herman/1988,
    13.7/Zuboff/2019
  } {
    \draw (\x, 0.15) -- (\x, -0.15);
    \node[name] at (\x, -0.25) {\t};
    \node[font=\tiny, anchor=south] at (\x, 0.2) {\d};
  }
\end{tikzpicture}
\end{center}
\vspace{1.5em}


\section{Notes and References}
\label{sec:ch3-notes}

On the sophists and rhetoric, see Plato's \textit{Gorgias} and \textit{Republic}.
\textcite{hobbes1651} and \textcite{mandeville1714} establish the selfish anthropology; on its selective amplification, see \textcite{mcCloskey2006}.
\textcite{marx1846} is the foundational text on ideology as material production; \textcite{althusser1970} develops the institutional apparatus.
\textcite{gramsci1971} develops hegemony and counter-hegemony.
\textcite{lippmann1922} and \textcite{dewey1927} frame the public-opinion debate.
\textcite{bernays1928} codifies propaganda technique; \textcite{ellul1965} extends it to a total social phenomenon.
\textcite{adorno1944} develops the culture industry concept.
\textcite{foucault1971} restructures the analysis of discourse and truth-production.
\textcite{bourdieu1984, bourdieu1991} provide the symbolic-power framework.
\textcite{herman1988} extends these ideas into a testable media model.
\textcite{said1978} and \textcite{fanon1961} analyze colonial knowledge production.
\textcite{zuboff2019} diagnoses platform-mediated surveillance capitalism.
On the institutional history of neoliberalism as a doxonomic project, see \textcite{mirowski2009}, \textcite{mayer2016}, and \textcite{slobodian2018}.
\textcite{stanley2015} provides a recent philosophical analysis of propaganda in democratic societies.
\textcite{hall1980} on encoding/decoding introduced the concept of negotiated readings that resist the dominant code.


\section{Exercises}

\begin{exercise}
For each thinker in this chapter, state the corresponding doxonomic concept from the formal framework of Chapter~\ref{ch:what-is-doxonomics} (e.g., Marx $\to$ endogenous $\fundingvec$; Bernays $\to$ applied optimization of $\Phi$).
\end{exercise}

\begin{exercise}
Bernays staged the ``Torches of Freedom'' campaign (1929) to associate women's smoking with feminist liberation.
Analyze this as a doxonomic operation: identify $\funding{k}$, $\credibility{k}$, $\relevance{k}{i}$, $\reach{k}$, and $\persistence{k}$.
Was the campaign operating at layer 1 (propositional belief) or at deeper layers?
\end{exercise}

\begin{exercise}
Compare Foucault's ``order of discourse'' with the concept of credibility weights $\credibility{k}$.
In what sense does Foucault's analysis go beyond the first-order model?
Write down a second-order model in which $\credibility{k}$ is itself produced by a meta-level doxonomic system, and discuss the implications for measurement.
\end{exercise}

\begin{exercise}
The propaganda model identifies five structural filters.
For each filter, write the corresponding formal constraint on the production function $\Phi$ or the narrative space $\narr$.
Which filters are most affected by the transition from legacy media to platform-based media?
\end{exercise}

\begin{exercise}
Choose a colonial education system (e.g., British India, French West Africa, Belgian Congo).
Reconstruct its belief production chain.
How did it differ from the neoliberal chain described in Example~\ref{ex:neoliberal-chain}?
What institutional features made colonial belief production particularly durable?
\end{exercise}

\begin{exercise}
The chapter presents the Lippmann-Dewey debate as mapping onto a tension between diagnostic and constructive doxonomics.
Is it possible to practice both simultaneously, or does the diagnostic stance (``belief is manufactured'') undermine the constructive stance (``democratic belief formation is possible'')?
Argue your position.
\end{exercise}

\begin{exercise}
Identify a doxonomic insight from a tradition \emph{not} covered in this chapter, for example from Islamic jurisprudence, Confucian statecraft, Buddhist epistemology, or indigenous knowledge systems.
State the insight in formal doxonomic language.
\end{exercise}
